%! Addition and Multiplication Rules — Unit 3, Lesson 3.3 — Homework (Answer Key)
% This is the lesson's SCORED (1 to 5) individual practice and the LAST component in the packet.
% ASSIGNED at Close & Assign and done OUTSIDE class. DELTAMATH + JUSTIFICATIONS form: items 1, 4,
% and 6 are turned in and keep their \gradedblank boxes (copied verbatim to the key), 2 points
% each = 6; DeltaMath carries items 2, 3, and 5 (plain \blank here, \ans in the key, an unboxed
% worked reference) and is scaled to the other 6. 12 points in all.
% THIRD CONTEXT: notes = concert raffle; Guided Practice = rain and the bus; here = an arcade.
% Items 1-2 on page 1, items 3-6 on page 2. Item 1 core + interpret; item 2 SPIRAL (Lesson 3.2); items 3-4 multiplication, with and
% without replacement (item 4 justify); item 5 the CONTRAST PAIR (blue/star has an overlap and
% is independent; blue/skull has none and is dependent); item 6 the CRUX head-on.
% THE DATA (verified in Python)
%   Wheel, 20 slices: B 8 (.4), S 5 (.25), both 2 (.1), B only 6, S only 3, outside 9 (includes
%     the 4 skulls K, .2).  B or S .55; not S .75; neither .45; P(B)P(S)=.1 -> independent;
%     P(B and K)=0, P(B)P(K)=.08 -> dependent; P(S and K)=0, P(S)P(K)=.05 -> dependent.
%   Tray, 5 capsules: 2 gold, 3 silver.  Put back: GG 4/25=.16; at least one 1-9/25=.64.
%     Kept: GG (2/5)(1/4)=.10; second gold after first gold .4 -> .25.
\documentclass[12pt]{article}
\usepackage{saar-article}
\usepackage{saar-key}   % pulls in -boxes; do NOT also load -boxes
\usepackage{grading}
\gradingassignment{saar-u03-l03-hw}
\gradingcourse{stats_analysis_algebraic_reasoning}
\newcommand{\TallMath}[1]{$\displaystyle #1\rule[-1.4em]{0pt}{3.2em}$}
% Handwriting room inside every \begin{work} block. Moves the blank and the
% key together, so the two can never drift apart.
\setlength{\workrowsep}{4pt}

\begin{document}

\pageheader{Unit 3, Lesson 3.3}{Homework\strut --- Answer Key}
{\small\itshape Write each final answer in its box. Only what is inside a box is graded.}\par
% NAMESTRIP: no \namedateperiod here — mirrors the blank. NO teachernote here —
% teacher prose lives in the lesson plan.

\vspace{0.08in}

% Authored identically in homework/ and homework_key/.
\begin{remindbox}
\small
\textbf{This is your graded homework.} It is scored 1 to 5 and due at the start of the first
class after two study halls. It is assigned today and done on your own, outside class --- flip
back to the \emph{Keep in Mind} box on the cover and to row 4 of your notes.
\end{remindbox}

\vspace{0.06in}

\boxguard[20]
\begin{scenariobox}[The Arcade Prize Games]{cerulean}
\small
\begin{minipage}[c]{0.52\linewidth}
\textbf{The wheel.} A prize wheel has $20$ equal slices, and one spin lands on one slice at
random. $B$: the slice is blue. $S$: the slice has a star. The Venn diagram counts slices. Four
of the slices outside both circles show a skull ($K$); a skull slice is never blue and never
has a star.
\end{minipage}\hfill
\begin{minipage}[c]{0.44\linewidth}\centering
\begin{tikzpicture}[x=1cm, y=1cm]
  \draw[thick, charcoal] (0,0) rectangle (6.8,2.75);
  \fill[cerulean, fill opacity=0.18] (2.6,1.3) circle[radius=1.1];
  \fill[goldacc, fill opacity=0.25] (4.1,1.3) circle[radius=1.1];
  \draw[thick, cerulean] (2.6,1.3) circle[radius=1.1];
  \draw[thick, goldacc] (4.1,1.3) circle[radius=1.1];
  \node[font=\footnotesize\bfseries] at (1.2,2.45) {$B$};
  \node[font=\footnotesize\bfseries] at (5.5,2.45) {$S$};
  \node[font=\large\bfseries] at (2.0,1.3) {6};
  \node[font=\large\bfseries] at (3.35,1.3) {2};
  \node[font=\large\bfseries] at (4.7,1.3) {3};
  \node[font=\large\bfseries] at (6.2,0.4) {9};
\end{tikzpicture}
\end{minipage}
\par\vspace{4pt}
\textbf{The capsule tray.} A tray holds $5$ capsules: $2$ gold and $3$ silver. A player picks one
capsule at random, then a second.
\end{scenariobox}

\vspace{0.04in}

\boxguard
\begin{notesbox}{Rules on the Wheel and the Tray}
Give each probability as a decimal.
\begin{enumerate}[leftmargin=1.6em, itemsep=5pt, topsep=2pt]

  \item Use the addition rule to find the probability of blue, star, or both on one spin.

        \begin{work}
          P(B \text{ or } S) &= 0.40 + 0.25 - 0.10
        \end{work}

        $P(B \text{ or } S) =$ \blankbox{\gradedblank[id=q1a, type=short, points=1, answer={0.55}, check=numeric, display={0.55}]{2.2cm}{13.5pt}}

        \gradedresponse[What does this number tell you about spins of this wheel?]{\gradedblank[id=q1b, type=rubric, points=1, display={In the long run, about 55\% of spins land on a blue slice, a star slice, or a blue star slice.}]{\linewidth}{2in}}

  \item \emph{Back to Lesson 3.2.} Read the Venn diagram.

        $P(\text{not } S) =$ \ans{0.75} \qquad
        $P(\text{neither } B \text{ nor } S) =$ \ans{0.45}

\end{enumerate}
\end{notesbox}

\boxguard[30]
\begin{notesbox}{Rules on the Wheel and the Tray, continued}
\begin{enumerate}[leftmargin=1.6em, itemsep=3pt, topsep=1pt, start=3]

  \item The player \textbf{puts the first capsule back} before picking again.

        \begin{work}
          P(\text{at least one gold}) &= 1 - (0.6)(0.6)
        \end{work}

        $P(\text{gold, then gold}) =$ \ans{0.16} \qquad
        $P(\text{at least one gold}) =$ \ans{0.64}


  \item Now the player \textbf{keeps the first capsule out} before picking again.

        $P(\text{gold, then gold}) =$ \blankbox{\gradedblank[id=q4a, type=short, points=1, answer={0.1}, check=numeric, display={0.10}]{2.2cm}{13.5pt}}

        \gradedresponse[Independent or dependent? Justify with the chance of gold on the second pick.]{\gradedblank[id=q4b, type=rubric, points=1, display={Dependent. Gold on the second pick is 0.4 before anything is taken, but only 0.25 after a gold is kept out, so the first pick changes the chance.}]{\linewidth}{2in}}

  \item Multiply, compare with the overlap, and write \emph{independent} or \emph{dependent}.

        \begin{work}
          P(B \text{ and } S) &= 0.10, \quad P(B \text{ and } K) = 0
        \end{work}

        $P(B) \cdot P(S) =$ \ans{0.10} \quad
        Blue and star: \ans{independent}

        \vspace{4pt}
        $P(B) \cdot P(K) =$ \ans{0.08} \quad
        Blue and skull: \ans{dependent}

  \item Tessa says, \emph{``A slice can't have a star and a skull, so those events must be
        independent.''} Is she right? Answer with numbers, then say what they mean for a spin.

        \gradedresponse{\gradedblank[id=q6, type=rubric, points=2, display={No. P(S and K) = 0, but P(S) times P(K) = (0.25)(0.2) = 0.05, so star and skull are dependent. Once you know the slice has a star, the chance of a skull drops from 0.2 to 0.}]{\linewidth}{2in}}

\end{enumerate}
\end{notesbox}

\boxguard[3]   % the short spiralbox needs about three baselines; a larger guard makes a third page
\begin{spiralbox}
\small
\textbf{Next class: Lesson 3.4, Conditional Probability, Independence, and Simulation.} Today's
second-branch numbers get a name and a notation, and we test independence with them.
\end{spiralbox}

\end{document}
